\section{Methods}\label{sec:methods}

We consider on-policy reinforcement learning for autoregressive policies with trajectory-level rewards. For an input prompt $x$ and sampled completion $y = (y_1, \ldots, y_T)$, the policy $\pi_\theta$ receives a scalar reward $R(x,y)$ from a verifier. The optimization target is
\begin{equation}
J(\theta) = \mathbb{E}_{x \sim \mathcal{D},\, y \sim \pi_\theta(\cdot \mid x)}[R(x,y)].
\end{equation}
We focus on the critic-free regime used in much recent RLVR work: the reward is sparse, generation is on-policy, and no learned value function is trained.

\subsection{Uniform Sequence-Level Policy Gradients}

For a fixed $(x,y)$, the REINFORCE estimator is
\begin{equation}
 g_{\mathrm{rf}}(x,y) = R(x,y) \sum_{t=1}^{T} \nabla_\theta \log \pi_\theta(y_t \mid y_{<t}, x).
\end{equation}
A baseline can be subtracted without changing the expectation when it is independent of the sampled token at position $t$:
\begin{equation}
 g_{\mathrm{base}}(x,y) = A(x,y) \sum_{t=1}^{T} \nabla_\theta \log \pi_\theta(y_t \mid y_{<t}, x),
\end{equation}
where $A(x,y)$ is the resulting sequence-level advantage.

We study two common prompt-level baselines. Given $K$ rollouts $\{y^{(i)}\}_{i=1}^{K}$ for the same prompt, with rewards $\{R^{(i)}\}_{i=1}^{K}$:
\begin{equation}
 b_{\mathrm{RLOO}}^{(i)}(x) = \frac{1}{K-1}\sum_{j \neq i} R^{(j)},
\end{equation}
and the corresponding advantage is $A^{(i)} = R^{(i)} - b_{\mathrm{RLOO}}^{(i)}$. For GRPO-style normalization we use
\begin{equation}
 A_{\mathrm{GRPO}}^{(i)}(x) = \frac{R^{(i)} - \mu_R(x)}{\sigma_R(x) + \varepsilon},
\end{equation}
where $\mu_R$ and $\sigma_R$ are the within-prompt reward mean and standard deviation. In both estimators, the scalar advantage is broadcast uniformly over tokens.

\subsection{Intrinsic Token Weighting}

We replace uniform broadcasting with token weights $w_t(x,y)$:
\begin{equation}
 g_w(x,y) = A(x,y) \sum_{t=1}^{T} w_t(x,y) \, \nabla_\theta \log \pi_\theta(y_t \mid y_{<t}, x),
\end{equation}
subject to
\begin{equation}
 w_t(x,y) \ge 0, \qquad \sum_{t=1}^{T} w_t(x,y) = 1.
\end{equation}
Uniform weighting is recovered by $w_t = 1/T$.

We emphasize an important point up front: unless the weights are constant with respect to the sampled trajectory, this estimator is generally \emph{biased} relative to the standard sequence-level policy gradient objective. Our goal is therefore not to preserve exact unbiasedness, but to study whether a lightweight, policy-internal credit-allocation prior can produce a favorable optimization trade-off. This framing is important both conceptually and experimentally: intrinsic token weighting changes \emph{where} the update mass is placed within a trajectory, and may thereby change the effective objective pursued by optimization.

We parameterize each weighting rule through a non-negative token salience score $\alpha_t(x,y)$ and normalize it within the completion:
\begin{equation}
 w_t(x,y) = \frac{\alpha_t(x,y) + \delta}{\sum_{k=1}^{T}(\alpha_k(x,y) + \delta)},
\end{equation}
with a small floor $\delta > 0$ to avoid degenerate all-zero trajectories. In implementation, the weights are treated as detached scalars multiplying the score-function terms, so the method introduces no second-order derivatives.

\paragraph{Uniform weighting.}
\begin{equation}
 \alpha_t^{\mathrm{uniform}} = 1.
\end{equation}
This recovers the standard critic-free baseline.

\paragraph{Surprisal weighting.}
We define token surprisal as
\begin{equation}
 \alpha_t^{\mathrm{surp}}(x,y) = -\log \pi_\theta(y_t \mid y_{<t}, x).
\end{equation}
This emphasizes sampled actions that the current policy regarded as comparatively unlikely.

\paragraph{Entropy-reduction weighting.}
Let
\begin{equation}
 H_t^{\mathrm{pre}}(x,y) = -\sum_v \pi_\theta(v \mid y_{<t}, x) \log \pi_\theta(v \mid y_{<t}, x)
\end{equation}
be the predictive entropy before sampling token $y_t$. We define the local entropy reduction as
\begin{equation}
 \alpha_t^{\mathrm{ent}}(x,y) = \max\{0, H_t^{\mathrm{pre}}(x,y) - H_{t+1}^{\mathrm{pre}}(x,y)\},
\end{equation}
with $\alpha_T^{\mathrm{ent}} = 0$. This emphasizes tokens after which the model's next-step distribution becomes more concentrated.

\subsection{Why Continuous Weighting?}

We focus on continuous weighting rather than hard token masking for two reasons. First, it is a smaller intervention: every token still receives some update mass, but semantically salient positions can receive more. Second, continuous weights can be paired directly with any existing critic-free policy-gradient implementation without changing the rollout structure or verifier API. In this sense, intrinsic weighting is a drop-in modifier of the loss rather than a replacement for the surrounding RLVR pipeline.

\subsection{Controlled Benchmark Instantiation}

To isolate credit assignment from the confounds of large-scale LLM training, we instantiate the method in a controlled benchmark implemented in NumPy. The policy is a small contextual autoregressive model with a learned prompt projection, previous-token embedding, position embedding, tanh hidden state, and linear output head. Each task generates trajectories of length five with a known decomposition into two low-semantic filler positions and three reward-relevant positions. This makes it possible to evaluate not only final task success but also the fraction of update mass placed on semantically important decisions.

The benchmark contains two tasks. In \texttt{arithmetic\_trace}, the prompt contains two integers and the target output includes a carry token, a parity token, and a coarse answer bucket. In \texttt{program\_trace}, the prompt contains a symbolic operation and a small value bucket, and the target output includes an operation token, an argument token, and an output token. Rewards are binary and exactly verifiable. The full implementation is provided in the repository under \texttt{experiment/}.
